% SGA 3，第 VII_B 篇，英文重建稿——第 1 轮。
% 范围：权威底本分篇第 91--92 页，第 5.2.1.C 节全文，止于第 5.2.2 节之前。
% 合并阅读本页码： 604--605.
% 再版印刷页码： 593--594.
% 本范围内编者注： 157.
% 权威底本： Polo--Gille Exp7B-13oct24.pdf,
% SHA-256 D303690E0195654BBE7C3DAA5D89EA23E24CD709A3575217234AC19AE8836694.
% 曾参考 Jacob Reinhold 的 CC BY 4.0 Markdown 作为比对谱系。
% 本组件没有图表；全部数学均为原生 LaTeX。

\paragraph{\SGARefLink{sga3:VIIB:5.2.1.C}{5.2.1.C}.}
\phantomsection
\label{sga3:VIIB:5.2.1.C}
\footnote{\label{sga3:VIIB:nde:157}编者在这里回到原文，并考虑到前面的增补而作了缩短。}
特别地，当 \(\omega\) 的滤过定常，即存在整数 \(n_0\) 使得
\[
  \omega_n=\omega_{n_0}
  \qquad
  \text{对 }n_0\leq n<+\infty.
\]
时，上述两个条件\textup{（\SGARefLink{sga3:VIIB:5.2.1.A}{5.2.1.A} 和 \SGARefLink{sga3:VIIB:5.2.1.B}{5.2.1.B}）}成立。此时，
\[
  k[[\omega]]
  \big/
  \bigl((x^{p^r})\bigr)_{x\in\omega_r}
\]
同构于完备化张量积
\[
\begin{aligned}
  &
  \frac{k[[\omega_1]]}
       {\bigl((x^p)\bigr)_{x\in\omega_1}}
  \widehat\otimes_k
  \frac{k[[\omega'_2]]}
       {\bigl((x^{p^2})\bigr)_{x\in\omega'_2}}
  \widehat\otimes_k\cdots
  \\
  &\hspace{3.2em}
  {}\widehat\otimes_k
  \frac{k[[\omega'_{n_0}]]}
       {\bigl((x^{p^{n_0}})\bigr)_{x\in\omega'_{n_0}}}
  \widehat\otimes_k
  k[[\omega'_\infty]],
\end{aligned}
\]
其中 \(\omega'_n\) 为 \(\omega_{n-1}\) 在 \(\omega_n\) 中的补空间，\(\omega'_\infty\) 为
\[
  \omega_\infty=\omega_{n_0}
\]
在 \(\omega\) 中的补空间。

在情形 \textup{\SGARefLink{sga3:vii:paragraph:5-1:equation:i}{(i)}} 中，对充分大的 \(r\) 有 \(\omega_r=\omega\)；在情形 \textup{\SGARefLink{sga3:vii:paragraph:5-1:equation:ii}{(ii)}} 中，\(\omega\) 为有限维；在情形 \textup{(iii)} 中，对每个 \(r\) 有
\[
  I_r=0
\]
所以在这些情形中，\(\omega\) 的滤过显然都定常。在最后一种情形中，\(B\) 同构于形式幂级数代数 \(k[[\omega]]\)。因此，结合下文 \SGARefLink{sga3:VIIB:5.2.2}{5.2.2}--\SGARefLink{sga3:VIIB:5.2.5}{5.2.5}，上述讨论即推出定理。
